\section{Full CPT cannot distinguish the grid pairs}
\label{eval:section}
\label{sec:cpt-lower-bound}

We now apply the support and counting results to an arbitrary sentence of
full \(\CPT\), with cardinality and unrestricted finite set rank.  The
remaining task is to describe its actual evaluation by a counting formula
of fixed variable width.  The formula may depend on the common input size;
its width will depend only on the sentence.  This is sufficient because
the supported HF domains agree on \emph{all} formulas of that width.

We use the bounded set-term semantics recalled in the preliminaries,
as specified in \cite[Section~2]{ls2023}.
For an iteration with bound \(p\), write
\(P(n)=\lfloor p(n)\rfloor\).  Its unbounded sequence starts at
\(a_0=\emptyset\) and applies the bounded semantics of its step term.
It returns the first stationary value \(a_j=a_{j+1}\) if
\(j\le P(n)\) and all states through index \(j\) satisfy
\(|\TC(a_i)|\le P(n)\); otherwise it returns \(\emptyset\).
Here and below the bounds are nonnegative at the input sizes considered.
After a stationary value all later states are the same, so this stopping
rule is equivalent to requiring the size bound throughout the sequence.
In particular, testing stationarity at index \(P(n)\) requires computing
the candidate \(a_{P(n)+1}\).

\subsection{An invariant family containing the entire evaluation}

For an HF object \(x\), put
\[
 \langle x\rangle=\{x\}\cup\TC(x).
\]
For an atom \(x\), this means \(\langle x\rangle=\{x\}\): atoms have no
members.  A finite family is membership-transitive if it contains all
members of each of its objects.

Fix a CPT sentence \(\psi\), including its polynomial annotation.
An \emph{activation} is a syntax occurrence together with the values
assigned to its free variables when that occurrence is evaluated.
Use the following deterministic evaluation prescription.  Evaluate all
arguments of a function and both children of every Boolean connective.
For a comprehension, evaluate its condition on every member of the
evaluated range, and its body on every member for which the condition is
true.  For an iteration, test successive candidates until stationarity,
the first excessive transitive size, or the iteration cutoff is
detected.  Include the candidate needed for the final stationarity test.
These conventions specify a family of activations without choosing an
order on any set.

Let \(\mathcal T_\psi(A)\) be the membership closure of all arguments and
term values occurring in these activations, together with all input atoms
and \(\emptyset\).  A rejected candidate and all objects constructed while
evaluating it are included.  Formula truth values need not be represented
by HF objects.

\begin{lemma}[Polynomial invariant evaluation family]
\label{eval:trace}
For every fixed CPT sentence \(\psi\), there is an integer \(q\ge1\)
such that, on every sufficiently large finite input \(A\) of size \(n\),
\[
 |\mathcal T_\psi(A)|\le n^q.
\]
The family \(\mathcal T_\psi(A)\) is membership-transitive and invariant
under \(\Aut(A)\).  It includes the first excessive candidate of each
activated iteration and, when needed, its candidate at index \(P(n)+1\).
\end{lemma}

\begin{proof}
We first prove a polynomial bound for an arbitrary fixed subexpression
evaluated on arguments of polynomial hereditary size.  More precisely,
if the sum of the sizes of the families \(\langle a\rangle\) for its
arguments is at most \(R\), syntax induction bounds both its number of
recursive activations and the total number of hereditary objects in
these activations by a polynomial in \(n+R+2\).  The polynomial may depend
on the subexpression and on the annotation.

Variables and constants satisfy this assertion directly.  Pair and Union
add or collect only polynomially many existing objects.  Unique returns
a member or \(\emptyset\).  Card returns the finite ordinal equal to the
number of members of its argument, or zero on an atom.  An ordinal \(k\)
and all its recursive members comprise \(k+1\) objects, so this operation
also satisfies the assertion.  Composition uses the child bounds a fixed
number of times.

For a comprehension, the induction hypothesis bounds the hereditary size
of its range by a polynomial in \(n+R+2\), hence bounds its number of
members.  Each condition is evaluated on the parameters and one such
member.  Each required body evaluation has arguments of the same
polynomial hereditary bound.  Applying the child estimates at this
enlarged bound and multiplying by the number of range members gives a
polynomial total.  Taking the set of body values adds one object and
does not increase this estimate beyond another polynomial.

For an iteration, every state supplied to the next step has passed the
guard and has hereditary size at most \(P(n)+1\).  The parameters retain
their original bound \(R\).  By the induction hypothesis, a step on these
arguments uses polynomially many hereditary objects.  There are at most
\(P(n)+1\) step evaluations before the stopping rule applies.  This
includes the first candidate violating the size guard: it is the output
of a child evaluation on the preceding, bounded state, so it still has
a polynomial hereditary bound, possibly larger than \(p(n)\).
It also includes the last candidate needed to compare
\(a_{P(n)+1}\) with \(a_{P(n)}\).  The empty default adds only a constant.
This argument applies to a nested iteration by the same syntax
induction; it does not assume that the step is iteration-free.
Boolean and atomic formula evaluations only combine these bounds.
At the sentence, \(R=0\), and a fixed integer exponent \(q\) absorbs the
resulting polynomial for all sufficiently large \(n\).

Every operation commutes with automorphisms of the input.  The full
comprehension range is used, the truth condition is invariant under the
corresponding action on parameters, and the iteration stopping tests use
only equality and cardinality.  Thus an automorphism sends every
activation to an activation at the same syntax occurrence and sends its
value to the corresponding value.  Starting with the empty assignment
makes the whole family invariant.  Closing it under membership preserves
invariance and gives membership transitivity by definition.
\end{proof}

\subsection{Exact descriptions at actual arguments}

The next lemma separates two requirements.  We need correctness at the
argument tuples actually used during evaluation.  At each such tuple,
however, uniqueness of the output must hold over the entire HF domain
under consideration. An extra proposed output could otherwise create a
false continuation of an iteration and prevent its empty-set default.

\begin{lemma}[Bounded-width description of a CPT evaluation]
\label{eval:description}
Fix a CPT sentence \(\psi\).  There are an integer \(m\ge1\) and, for
each sufficiently large input size \(n\), a \(C^m\) sentence
\(\widehat\psi_n\) with the following property.

Let \(A\) be an \(n\)-element input and let \(D\subseteq\HF(A)\) be
membership-transitive, contain all atoms and all pure HF sets, and contain
\(\mathcal T_\psi(A)\).  Interpret membership and the atom predicate on
\(D\), and interpret each input relation as false if an argument is a
set.  Then
\[
 D\models\widehat\psi_n
       \quad\Longleftrightarrow\quad
 \psi\text{ evaluates to true on }A.
\]
The formula \(\widehat\psi_n\) depends only on \(\psi\) and \(n\), not on
\(A\) or \(D\).  Its length and counting constants may depend on \(n\);
\(m\) does not.
\end{lemma}

\begin{proof}
Choose a common integer \(J=J(n)\ge1\), depending only on the polynomial
bounds in the proof of Lemma~\ref{eval:trace}, which bounds the hereditary
size of every argument and value in the evaluation.  In particular it
bounds all cardinalities passed to Card and all descending membership
depths of tested candidates.  Increasing \(J\) to an integer polynomial
bound is harmless.

For each term occurrence \(t\), with free argument tuple \(\bar x\),
we construct a relational formula \(G_t(\bar x,v)\).
For each formula occurrence \(\theta\), we construct a truth formula
\(F_\theta(\bar x)\).  Write \(\operatorname{val}_A(t,\bar a)\) for the
bounded semantic value of \(t\) at \(\bar a\).
Our simultaneous induction assertion is
\begin{equation}\label{eval:graph-invariant}
 \begin{split}
 D\models G_t(\bar a,v)
   &\quad\Longleftrightarrow\quad
       v=\operatorname{val}_A(t,\bar a)
       \qquad\text{for every }v\in D,\\
 D\models F_\theta(\bar a)
   &\quad\Longleftrightarrow\quad
       \theta\text{ is true at }\bar a,
 \end{split}
\end{equation}
whenever the displayed argument tuple is an activation of the indicated
occurrence.  Every genuine term value belongs to \(D\) by the trace
assumption.  No assertion is made about an argument tuple which never
occurs.

Write \(\operatorname{At}(v)\) for the atom predicate and define the
typed empty-set predicate
\[
 E(v):=\neg\operatorname{At}(v)\ \land\
                    \forall z\,\neg(z\in v).
\]
It has exactly the empty set as its solution in \(D\).  Indeed a nonempty
set in \(D\) has a member in \(D\) by membership transitivity, while the
non-atom conjunct excludes urelements.

\paragraph{Ordinals and the elementary operations.}
There are two-variable formulas \(O_k(v)\) defining the pure von Neumann
ordinal \(k\), for every fixed nonnegative integer \(k\):
\[
 O_0(v):=E(v),\qquad
 O_k(v):=\neg\operatorname{At}(v)\ \land\
       \forall z\left(z\in v\ \longleftrightarrow\
                         \bigvee_{0\le j<k}O_j(z)\right)\quad(k>0).
\]
All pure ordinals belong to \(D\).  Induction on \(k\), membership
transitivity and extensionality show that \(O_k\) has exactly the
claimed solution.  The two names alternate in the recursion:
inside \(O_j(z)\), the old name \(v\) can be bound, since that child
formula has only \(z\) free.

Variable graphs are equalities with their corresponding arguments.
The graphs for \(\emptyset\) and \(\Atoms\) are respectively \(E(v)\) and
\[
 \neg\operatorname{At}(v)\ \land\
          \forall z\,(z\in v\longleftrightarrow\operatorname{At}(z)).
\]
For evaluated arguments \(a,b\), the graphs of Pair and Union are
\[
 \begin{split}
 &\neg\operatorname{At}(v)\ \land\
               \forall z\,(z\in v\longleftrightarrow(z=a\lor z=b)),\\
 &\neg\operatorname{At}(v)\ \land\
               \forall z\,(z\in v\longleftrightarrow
                            \exists w\,(w\in a\land z\in w)).
 \end{split}
\]
Put
\[
 \operatorname{Sing}(a,w):=
       \neg\operatorname{At}(a)\ \land\
                         \forall z\,(z\in a\longleftrightarrow z=w).
\]
The graph of Unique is
\[
 \operatorname{Sing}(a,v)\ \lor\
             \bigl(\neg\exists w\,\operatorname{Sing}(a,w)\land E(v)\bigr).
\]
These are graphs of the operations, not assertions that their outputs
are atoms or that their inputs are sets.  They give the specified
defaults on atoms as well.

The graph of Card on the actual argument \(a\) is
\[
 \bigl(\operatorname{At}(a)\land E(v)\bigr)\ \lor\
 \left(\neg\operatorname{At}(a)\land
   \bigvee_{k=0}^{J}
      \left((\exists^{=k}z\,z\in a)\land O_k(v)\right)\right).
\]
Every member of \(a\) belongs to \(D\), so the quantifier counts its
actual members.  Its actual cardinality is at most \(J\), and therefore
one and only one disjunct yields an output.

For a composite term, existentially bind the outputs of its child terms
and apply the appropriate operation graph.  By the induction hypothesis,
these intermediate outputs are forced to be their genuine values among
all elements of \(D\).  Membership transitivity then makes each displayed
extensional graph exact among all possible proposed outputs.  The graph
for an input relation applied to terms is obtained in the same way;
for equality use equality of their outputs.  Translate Boolean
connectives by the corresponding Boolean connectives.

\paragraph{Comprehension.}
Suppose
\[
 t(\bar x)=\{u(\bar x,y):y\in r(\bar x),\ \theta(\bar x,y)\}.
\]
Its graph is
\[
 \begin{split}
 G_t(\bar x,v):={}&
 \exists w\biggl(
 G_r(\bar x,w)\land\neg\operatorname{At}(v)\land\\[-2mm]
 &\hspace{7mm}
 \forall z\left[
 z\in v\ \longleftrightarrow\
 \exists y\bigl(y\in w\land F_\theta(\bar x,y)
                       \land G_u(\bar x,y,z)\bigr)
 \right]\biggr).
 \end{split}
\]
At an actual parameter tuple, \(w\) must be the actual range.  Each of its
members activates the condition.  When the condition is true, that
member also activates the body, whose graph has exactly its correct
output by induction.  When the condition is false, the conjunction is
false without any correctness requirement on the body graph at that
unused argument tuple.  This proves that the right-hand side of the
membership biconditional defines exactly the actual collected values.
An alternative output set with an extra member is also excluded:
that member lies in \(D\) and is tested by the universal quantifier.
The non-atom condition excludes an atom as a spurious empty output.

\paragraph{The transitive-size guard.}
For positive integers \(\ell\), define strict downward path formulas
\[
 \begin{split}
 R_1(r,z)&:=(z\in r),\\
 R_{\ell+1}(r,z)&:=
       \exists w\,(w\in r\land R_\ell(w,z)).
 \end{split}
\]
Three variable names suffice for all \(\ell\): the endpoints and one
intermediate name, with the former root name recycled inside the child
formula.  Set
\[
 R_{\le J+1}(r,z):=\bigvee_{\ell=1}^{J+1}R_\ell(r,z).
\]
For every actual value \(a\) under consideration, its solutions in \(D\)
are exactly \(\TC(a)\).  All descendants belong to \(D\), and the chosen
depth bound reaches every descendant.  The paths have positive length;
they do not count the root merely by a path of length zero.  Exact
counting counts endpoints, not paths, even if a descendant is reached
in several ways.

For an iteration bound \(p\), the formula
\[
 \operatorname{Ok}_{p,n}(r):=
       \bigvee_{k=0}^{P(n)}
               \exists^{=k}z\,R_{\le J+1}(r,z)
\]
therefore says exactly \(|\TC(r)|\le P(n)\) on every actual candidate.
It also handles atoms and \(\emptyset\), whose strict transitive closures
are empty.  Its behavior on unrelated objects of greater depth is not
needed.

\paragraph{Iteration and the empty default.}
Consider a term \(t=s[y]^*\), with free parameters \(\bar x\), and let
\(S(\bar x,w,v)\) be the already constructed graph for its step term.
At the present input size abbreviate \(P=P(n)\) and
\(\operatorname{Ok}=\operatorname{Ok}_{p,n}\).  Define finite formulas
\[
 \begin{split}
 I_0(\bar x,v)&:=E(v)\land\operatorname{Ok}(v),\\
 I_{j+1}(\bar x,v)&:=
       \operatorname{Ok}(v)\land
       \exists w\,(I_j(\bar x,w)\land S(\bar x,w,v)),
                  \qquad 0\le j<P.
 \end{split}
\]
Let
\[
 \begin{split}
 F(\bar x,v)&:=
       \bigvee_{j=0}^{P}
                   \bigl(I_j(\bar x,v)\land S(\bar x,v,v)\bigr),\\
 G_t(\bar x,v)&:=
       F(\bar x,v)\ \lor\
        \bigl(E(v)\land\neg\exists w\,F(\bar x,w)\bigr).
 \end{split}
\]
For \(P=0\) the disjunction has its \(j=0\) term: a single step still
tests whether the initial empty state is stationary.

We verify this graph for an actual parameter tuple \(\bar a\).
Induction on \(j\) shows that before stopping, \(I_j(\bar a,v)\)
has exactly the genuine state \(a_j\) as its solution if all states
through that index pass the size guard.  At an admitted state, the step
activation is actual.  Its graph has exactly its actual candidate as
output over the whole of \(D\), by the syntax induction hypothesis.
This candidate is in \(D\), including when it is the first excessive
one.  The exact guard either admits that candidate or makes
\(I_{j+1}\) have no solution.  In the latter case all subsequent \(I\)'s
have no solution, so no unrelated state can enter later.

If an allowed stationary state is reached early, all subsequent \(I_j\)
repeat that same state.  The step graph is still exact there because this
argument tuple was already an actual activation.  If no earlier stop
occurs, the test \(S(\bar a,a_P,a_P)\) uses the actual final candidate
\(a_{P+1}\), which was included in the evaluation family.  It succeeds
exactly when stationarity is reached at index \(P\).
Thus \(F\) is empty when the bounded iteration fails, and otherwise its
only output is the genuine stationary value.  The final displayed
formula returns precisely that value or the typed empty default.
This proves the term part of~\eqref{eval:graph-invariant}.
In particular, a fixed point at some argument tuple that was never
activated cannot supply an extra output: every successful disjunct
first requires a state reached from \(I_0\).

\paragraph{One fixed supply of variable names.}
All formulas above are finite at fixed \(n\).  It remains to show that
their variable width is bounded independently of \(n\).
For each syntax occurrence, reserve its finitely many argument names,
an output name when it is a term, and finitely many names for the
intermediate outputs and local quantifiers in its displayed construction.
Child formulas receive their indicated free interface and a separate
working supply, renamed to avoid capture.  Since the syntax tree is
fixed, the sum of these finite reservations depends only on \(\psi\).

The only repetitions whose depths grow with \(n\) are the ordinal,
membership-path and iteration expansions.  The first two use the
recycled interfaces already described.  In the iteration recursion,
retain \(\bar x\) and two state names.  When \(I_j(\bar x,w)\) is inserted
in the formula for \(I_{j+1}(\bar x,v)\), its internal quantifier may bind
the name \(v\): the later state is not free in that child.  This binder
has scope only inside \(I_j\), so it does not capture \(v\) in the
sibling formula \(S(\bar x,w,v)\).  The step graph uses a separate fixed
working supply.  Repeating the recursion therefore does not keep a
tuple of all previous states.  Identified slots such as
\(S(\bar x,v,v)\) can, if needed, use a fresh local alias constrained
equal to \(v\).

This argument applies at every syntax level, including inside step
terms that themselves contain iterations.  Nesting depth adds only a
fixed number of reservations.  Let \(m\) bound their total number.
Taking the translated truth formula at the top sentence gives
\(\widehat\psi_n\), and~\eqref{eval:graph-invariant} at the empty
assignment proves the claim.
\end{proof}

The descriptions are a tool for comparing two evaluations.  They need
not be computable by the CPT sentence, have polynomial length, or
describe evaluation on every assignment in \(D\).  Their relevant
properties are the fixed width and the exactness asserted in
Lemma~\ref{eval:description}.  Neither property puts a bound on the
finite ranks occurring in \(D\).

\subsection{Applying support and transfer}

Recall the two grid structures \(A_{b^0}\) and \(A_{b^1}\), of common
size \(N=N_L=\Theta(L^3)\). The group \(H\) acts on each original input,
while its central subgroup \(K\) also preserves the analytical expansion
of Section~\ref{sec:base-counting}. We now combine these two actions:
the first places all actual evaluation values in a supported domain,
and the second lets us compare those domains by counting logic.

\begin{theorem}[Eventual equality for every CPT sentence]
\label{eval:pair-equivalence}
For every sentence \(\psi\) of full CPT with counting over \(\tau\),
there is \(L_\psi\) such that for every \(L\ge L_\psi\),
\[
 A_{b^0}\models\psi
       \quad\Longleftrightarrow\quad A_{b^1}\models\psi .
\]
\end{theorem}

\begin{proof}
Fix \(\psi\) first.  Choose the integer \(q\) from
Lemma~\ref{eval:trace} and the width \(m\) from
Lemma~\ref{eval:description}.  For each \(L\), put
\[
 s=\left\lceil2L(q\log_3N+1)^2\right\rceil+1,
 \qquad M=\lfloor L^{7/4}\rfloor.
\]
For \(i\in\{0,1\}\), let \(D_{L,i}^{(s)}\) consist of the objects of
\(\HF(A_{b^i})\) such that the object and each recursive member have
a \(K\)-support of length at most \(s\).
Give this domain membership, the atom predicate, and the original input
relations, false whenever any argument is not an atom.

The family \(\mathcal T_\psi(A_{b^i})\) is \(H\)-invariant,
membership-transitive and of size at most \(N^q\), for large \(L\).
Corollary~\ref{sup:hereditary} therefore gives
\[
          \mathcal T_\psi(A_{b^i})\subseteq D_{L,i}^{(s)}.
\]
The domain is membership-transitive and contains all atoms and pure HF
sets: an atom supports itself, and a pure set is fixed by every atom
permutation.  Thus it satisfies all the domain requirements of
Lemma~\ref{eval:description}.

For the fixed \(q,m\), Corollary~\ref{base:width} supplies all the
quantitative hypotheses of Theorem~\ref{hf:transfer} on sufficiently
large grids: the two analytical expansions are \(C^M\)-equivalent,
equal \(C^M\) types of tuples of length at most \(\max\{2,m\}s\)
determine \(K\)-orbits, and \(\max\{4,m+1\}s\le M\).
These are nonempty finite structures in the same binary vocabulary,
and \(K\) acts by automorphisms of both expansions.

The supported domains computed from the expansions have the same
underlying sets \(D_{L,i}^{(s)}\), since supports depend on the atom action
of \(K\), which did not change.  The theorem gives \(C^m\) equivalence
with the additional analytical relations; forgetting those relations
gives
\[
                 D_{L,0}^{(s)}\equiv_{C^m}D_{L,1}^{(s)}.
\]
The larger group \(H\) was used only on the original input and its
evaluation family.  No assertion that \(H\) preserves the analytical
expansion is needed.

Both inputs have size \(N\), so the same formula
\(\widehat\psi_N\), with the same numerical cutoffs and ordinal
predicates, describes their evaluations.
Lemma~\ref{eval:description} and the displayed \(C^m\) equivalence
therefore give equal truth values for \(\psi\) on the two inputs.
All inequalities hold eventually in \(L\) after fixing \(\psi,q,m\),
which proves the quantified assertion.
\end{proof}

\begin{proof}[Proof of Theorem~\ref{thm:main}]
Let \(\Phi\) be the single sentence constructed in
Section~\ref{sec:sentence}, over the eight-relation vocabulary \(\tau\)
and with its fixed polynomial annotation. It has one syntactic WSC
occurrence, and Proposition~\ref{sen:all-input} gives a Boolean value
on every finite \(\tau\)-structure. Let \(Q\) be its true-model class.
By Theorem~\ref{grid:separation}, for every \(L\ge2\),
\[
              A_{b^0}\in Q,\qquad A_{b^1}\notin Q.
\]
If a CPT sentence \(\psi\) defined \(Q\), it would have opposite truth
values on every such pair.  This contradicts
Theorem~\ref{eval:pair-equivalence} for large \(L\).
Finally, the non-strict inclusion follows from the inclusion of CPT
syntax in CPT+WSC. Thus
\[
       \operatorname{Def}_{\CPT}(\tau)
         \ \subsetneq\
       \operatorname{Def}_{\CPT+\WSC}(\tau).
\]
\end{proof}
